\documentclass[landscape]{beamer}
\usetheme{default}
\usepackage{amsmath, amssymb}
\setbeamertemplate{navigation symbols}{}

\title{4.3 Similarity \\ Linear Algebra Made Human}
\author{MATH 304 Rewritten for Clarity}
\date{}

\begin{document}

%------------------------------------------------
\begin{frame}
\titlepage
\centering
\Large
Take a breath.

Different bases do NOT change the transformation.

They only change the camera angle.

Today: We learn how matrices change when the basis changes.

And why it is actually simple.
\end{frame}

%------------------------------------------------
\begin{frame}{Cheat Sheet (1/4)}
\begin{columns}

\column{0.5\textwidth}
\textbf{Linear Operator}

$L: V \to V$

\textbf{Meaning:}
\begin{itemize}
\item $V$ = vector space
\item $L$ takes vectors in $V$
\item Outputs vectors in same $V$
\end{itemize}

\textbf{Key Idea:}
Matrix depends on basis.

\column{0.5\textwidth}

\textbf{Matrix Representation}

If $E=\{v_1,\dots,v_n\}$

Then matrix $A$ satisfies:

\[
A[x]_E = [L(x)]_E
\]

\textbf{Symbols:}
\begin{itemize}
\item $[x]_E$ = coordinate vector
\item $A$ = matrix in basis $E$
\end{itemize}

\end{columns}
\end{frame}

%------------------------------------------------
\begin{frame}{Cheat Sheet (2/4)}
\begin{columns}

\column{0.5\textwidth}
\textbf{Transition Matrix}

Let $E=\{v_i\}$

Let $F=\{w_i\}$

$S$ = transition matrix
from $F$ to $E$

\[
[x]_E = S[x]_F
\]

\textbf{Purpose of $S$:}

Changes coordinate systems.

\column{0.5\textwidth}

\textbf{Inverse Transition}

\[
[x]_F = S^{-1}[x]_E
\]

If $S$ did not exist:

We could not translate
between coordinate systems.

It is the translator.

\end{columns}
\end{frame}

%------------------------------------------------
\begin{frame}{Cheat Sheet (3/4)}
\begin{columns}

\column{0.5\textwidth}
\textbf{Similarity Formula}

If:
\begin{itemize}
\item $A$ = matrix of $L$ in $E$
\item $B$ = matrix of $L$ in $F$
\item $S$ = transition $F \to E$
\end{itemize}

Then:

\[
B = S^{-1} A S
\]

\column{0.5\textwidth}

\textbf{Meaning of each piece}

$S$ = move from $F$ to $E$

$A$ = apply transformation

$S^{-1}$ = move back to $F$

Without $S^{-1}$:

Result would stay in wrong coordinates.

\end{columns}
\end{frame}

%------------------------------------------------
\begin{frame}{Cheat Sheet (4/4)}
\begin{columns}

\column{0.5\textwidth}
\textbf{Definition: Similar Matrices}

$A$ and $B$ are similar if

\[
B = S^{-1} A S
\]

for some invertible $S$.

\textbf{Invertible =}
$\det(S) \neq 0$

\column{0.5\textwidth}

\textbf{Core Fact}

If $A$ and $B$ represent
the same operator
in different bases,

then they are similar.

Similarity = same transformation,
different viewpoint.

\end{columns}
\end{frame}

%------------------------------------------------
\begin{frame}{Big Picture Intuition}
Think basketball replay.

Same play.

Different camera angle.

The play did NOT change.

Only perspective changed.

Matrix $A$ and $B$:

Same transformation.

Different coordinate lens.

Similarity = changing cameras.
\end{frame}

%------------------------------------------------
\begin{frame}{Theorem (Unpacked)}
Let:

\begin{itemize}
\item $E=\{v_1,\dots,v_n\}$
\item $F=\{w_1,\dots,w_n\}$
\item $L:V\to V$
\item $S$ = transition from $F$ to $E$
\item $A$ = matrix of $L$ in $E$
\item $B$ = matrix of $L$ in $F$
\end{itemize}

Then:

\[
B = S^{-1}AS
\]

\end{frame}

%------------------------------------------------
\begin{frame}{Why Each Symbol Exists}
$S$:

Converts $F$-coordinates to $E$.

If removed:
$A$ cannot act correctly.

$A$:

Applies transformation in $E$.

If removed:
No transformation happens.

$S^{-1}$:

Returns answer to $F$ basis.

If removed:
Output stuck in wrong system.

Each symbol has a job.
\end{frame}

%------------------------------------------------
\begin{frame}{Example: $\mathbb{R}^2$ (Part 1)}
Given:

\[
A=\begin{pmatrix}
2 & 0\\
1 & 1
\end{pmatrix}
\]

New basis:

\[
u_1=\begin{pmatrix}1\\1\end{pmatrix},\;
u_2=\begin{pmatrix}-1\\1\end{pmatrix}
\]

Transition matrix:

\[
U=(u_1\;u_2)
=
\begin{pmatrix}
1 & -1\\
1 & 1
\end{pmatrix}
\]

\end{frame}

%------------------------------------------------
\begin{frame}{Example (Part 2)}
Compute determinant:

\[
\det(U)=1(1)-(-1)(1)=2
\]

So invertible.

\[
U^{-1}
=
\frac{1}{2}
\begin{pmatrix}
1 & 1\\
-1 & 1
\end{pmatrix}
\]

We can now compute:

\[
B=U^{-1}AU
\]

\end{frame}

%------------------------------------------------
\begin{frame}{What Just Happened?}
Step 1: Convert new basis to standard.

Step 2: Apply $A$.

Step 3: Convert back.

That is literally it.

It looks scary.

But it is just
translate → act → translate back.

Like switching game engines,
running physics,
then switching back.
\end{frame}

%------------------------------------------------
\begin{frame}{Differentiation Example}
Operator: $D$

$D(p)=p'$

Basis 1:

$[1,x,x^2]$

Matrix:

\[
B=
\begin{pmatrix}
0&1&0\\
0&0&2\\
0&0&0
\end{pmatrix}
\]

Why?

$D(1)=0$

$D(x)=1$

$D(x^2)=2x$

\end{frame}

%------------------------------------------------
\begin{frame}{Same Operator, New Basis}
New basis:

$[1,2x,4x^2-2]$

Matrix becomes:

\[
A=
\begin{pmatrix}
0&2&0\\
0&0&4\\
0&0&0
\end{pmatrix}
\]

Different matrix.

Same differentiation.

Just different coordinate scaling.

Same TikTok.
Different filter.
\end{frame}

%------------------------------------------------
\begin{frame}{Homework Survival Rule}
If homework says:

"Find matrix in new basis"

You ALWAYS:

1. Build transition matrix.
2. Compute inverse.
3. Use $S^{-1}AS$.

Or:

Express $L(w_i)$
as combinations of basis vectors.

Two methods.
Same result.
\end{frame}

%------------------------------------------------
\begin{frame}{Extra Example (Deep WHY) Part 1}
Suppose:

\[
A=
\begin{pmatrix}
3&0\\
0&5
\end{pmatrix}
\]

Diagonal matrix.

Why is this nice?

It stretches independently.

Like scaling x and y separately.
\end{frame}

%------------------------------------------------
\begin{frame}{Extra Example (Part 2)}
Pick invertible

\[
S=
\begin{pmatrix}
1&1\\
0&1
\end{pmatrix}
\]

Compute:

\[
B=S^{-1}AS
\]

Result:

Matrix no longer diagonal.

But transformation
still stretches in two directions.

Similarity hides structure.
Does not destroy it.
\end{frame}

%------------------------------------------------
\begin{frame}{Final Takeaway}
Similarity means:

Same linear operator.

Different coordinate system.

Nothing mystical.

Just camera angles.

Translate → Transform → Translate Back.

You are not bad at this.

It was just presented without the story.
\end{frame}

\end{document}